% Documentclass Settings
	\documentclass[a4paper, twoside]{article}
	% Margins
	\usepackage{geometry}
		\geometry{
			% showframe,
			left   = 20mm,
			right  = 20mm,
			top    = 20mm,
			bottom = 20mm
		}

	% Headers and footers
	\usepackage{fancyhdr}
		\pagestyle{fancy}
		\fancyhf{}
		\fancyhead[LE, RO]{\thepage}
		\fancyhead[CE]{\textsc{Cong Wen}}
		\fancyhead[CO]{\textsc\rightmark}

		\setlength{\headheight}{15pt}
		\renewcommand\headrulewidth{0pt}








\input{../universal-pre}

\title{\tbf{ABS2017}}
\author{\tbf{Notes} by Cong Wen, 12/2022}
\date{}
% Last Updated: July 11, 2022



\begin{document}
\maketitle
\thispagestyle{empty}
\begin{abstract}
	This is my self-use notes\footnote{Last updated on January 19, 2022} on \tbf{On hyperplane sections of K3 surfaces}\cite{ABS2017}. Please reach out if you find anything incorrect.

\end{abstract}
\tableofcontents
% ----------------------------------------------------------------------------------------------------

\sct{The vanishing theorem}

The paper generalized a theorem of Wahl\cite{Wahl1997}, characterizing certain BNP curves in terms of Wahl map. Precisely stated as below:

\begin{thm}[Theorem~1.1]
	Let $C$ be a BNP curve of genus $g\gqs12$, then $C$ lies on a polarized K3 surface (or a limit of which) if and only if its Wahl map is not surjective.
\end{thm}

To learn the proof, it contains two key parts, the first part proves a vanishing result using $g\gqs12$, the second part uses BNP condition to complete the geometric argument.

\sct{Vanishing of the square of an ideal sheaf}

\begin{thm}[Theorem~1.3]
	Let $C\sbs\bP^{g-1}$ be a canonically embedded curve of $g\gqs11$ and $\mathrm{Cliff}(C)\gqs3$, then
	\[
		\mathrm{H}_{}^{1}(\bP^{g-1},\cI_{C/\bP^{g-1}}^{2}(k))=0,\fal k\gqs3.
	\]
\end{thm}


First, recall a theorem of Voisin \cite{Voisin1988},
\begin{thm}[Theorem~2.3]
	If $\mathrm{Cliff}(C)\gqs3$, $L$ a general element of a component $S$ of $W_{g-2}^{1}(C)$, then
	\begin{enr}
		\item $L$ and $K_{C}-L$ are both bpf.
		\item $\phi_{\sP{K_{C}-L}}$ maps $C$ birationally onto $\Gma\sbe\bP^{2}$ of degree $g$ with only double points as singularities.
		\item $L$ is BNP.
	\end{enr}
\end{thm}

Then, identify $C$ with its canonically embedded curve in $\bP:=\bP^{g-1}$, take some $L$ as above, construct a cubic scroll $C\sbs X\sbs\bP^{g-1}$.

Now, geometrically, $X$ is a cone over the Segre variety $\bP^{1}\tms\bP^{2}\sbs\bP^{5}$ with a $(g-7)$-dimensional vertex $\bP W$.


Then the followings are true:

\begin{lem}
	\begin{enr}
		\item All quadrics $Q\in \mathrm{H}_{}^{0}(\bP^{g-1},\cI_{X/\bP^{g-1}}(2))$ have rank at most $4$.
		\item Any $p\in X\bs\bP W$, there is a unique up to a constant factor non-zero quadric $Q_{p}\in \mathrm{H}_{}^{0}(\bP,\cI_{X/\bP}(2))$ which is singular at $p$.
		\item If a quadric $Q\in \mathrm{H}_{}^{0}(\bP,\cO_{\bP}(2))$ contains $R\cup C$ for some ruling $R\sbs X$ and is singular at a point $p\in C$ with $p\nin R$, then $Q=Q_{p}\in \mathrm{H}_{}^{0}(\bP,\cI_{X/\bP}(2))$.
	\end{enr}
\end{lem}




The key part of the paper goes as follows,
\begin{thm}
	If $C$ is a BNP curve of genus $g\gqs12$, then $C$ lies on a polarized K3 surface or on a limit thereof if and only if its Wahl map is not surjective.
\end{thm}

\begin{prp}
	Let $\oL{S}\sbs\bP^{g}$ be a surface with canonical sections, and $C$ be a general hyperplane section of $\oL{S}$, assume $C$ is BNP of genus $g\gqs12$. Then the following are true.
	\begin{enr}
		\item If $S_{0}=\bP(E)$ is ruled over a smooth curve $\Gma$ of genus $\gma$, then $\gma\leq 1$.
		\begin{enr}
			\item If $S_{0}=\bP(E)$ is ruled over an elliptic curve $\Gma$, then $e+1\leq h\leq 7$.
			\item If $S_{0}=\bP(E)$ is ruled over a smooth rational curve $\Gma$, then this minimal model can be replaced by $\bP^{2}$.
		\end{enr}
		\item If $S_{0}=\bP^{2}$, then $10\leq h\leq 18$, and $K_{S}^{2}=9-h\gqs-9$.
	\end{enr}
\end{prp}


One key proposition involved in the computation comes from Epema\cite{Epema1984}.

\begin{prp}[Proposition~9.3]
	Suppose $S_{0}=\bP(E)$, by replacing $S_{0}$ with another minimal model, the following hold
	\begin{enr}
		\item $E=\cO_{\Gma}\op\cO_{\Gma}(D), \Deg D=-e$.
		\item Either $D\sim-K_{\Gma}$, or $h\gqs1$ and $e>2\gma-2$.
		\item $\nu_{i}\leq a-1, i=1,\ldots,h$.
		\item $\sum\nu_{i}=a(e-2\gma+2)$.
		\item $2g-2=a^{2}e-\sum\nu_{i}^{2}$.
		\item $2a^{2}(\gma-1)+a(e-2\gma+2)\leq2g-2\leq a^{2}e-a(e-2\gma+2)$.
		\item $(a-1)h\leq2[g-1-a^{2}(\gma-1)]$ for $a\gqs2$.
		\item None of the base points of $\sP{C_{0}}$ lie on $\sgm\sbs S_{0}$.
	\end{enr}
\end{prp}

\sct{The geometric argument}

\begin{stp}
	A surface with canonical sections is a projective surface $\oL{S}\sbs\bP^{g}$ with $g\gqs3$ whose general section is a smooth canonical curve of genus $g$. Such surfaces are first studied and classified by Epema, here it is required to be not a cone. Then we have the diagram
	\[
		\begin{tikzcd}
			S_{0} & S\ar[r, "q"]\ar[l, "p"] & \oL{S}\sbs\bP^{g},
		\end{tikzcd}
	\]
	where $q$ is the minimal resolution and $S_{0}$ is a choice of a minimal model of $S$. $S$ is smooth and equipped with a big and nef divisor $C$.
\end{stp}

In such a case, the strucure theorem is given in Epema's book\cite{Epema1984}:
\begin{prp}
	\begin{enr}
		\item $\oL{S}$ is projectively normal with $h^{0}(-K_{S})=1$.
		\item $S$ is either a minimal K3 surface or a ruled surface.
		\item If $\oL{S}$ is not K3, then $S_{0}$ is either $\bP^{2}$ or a $\bP^{1}$-bundle over a smooth genus $\gma$ curve $\Gma$.
	\end{enr}
\end{prp}


\begin{thm}[Theorem~10.3]
	If $\oL{S}$ has smoothable singularities then $[\oL{S}]\in\Hilb^{\bP^{g}}$ belongs to $\oL{\cM}$, where $\oL{\cM}$ is the family of non-singular K3 surfaces of degree $2g-2$.
\end{thm}



% ----------------------------------------------------------------------------------------------------
\bibliographystyle{alpha}
\bibliography{../universal-ref}
\end{document}

